%!TEX TS-program = pdflatex
%!TEX TS-options = -shell-escape
% % % % %   Die folgenden Zeilen müssen ihre Zeilennummern 4 und 5 behalten !!!    % % % % %
\newcommand{\printpraesenzlsg}{false}
\newcommand{\printloesungen}{true}
\newcommand{\printbewertungen}{false}
% % % % %   \newcommand{\printloesungen}{false}                                    % % % % %
\newcommand{\blattnummer}{1}
%\newcommand{\abgabetermin}{\textcolor{red}{bis 11.04.2022, 09:00 Uhr}}
\documentclass[%
  paper=a4,
  fontsize=10pt,
  ngerman
  ]{scrartcl}

% Basics für Codierung und Sprache
% ===========================================================
  \usepackage{scrtime}
  \usepackage{etex}
  \usepackage{shellesc}
  \usepackage[final]{graphicx}
  \usepackage[utf8]{inputenc}
  \usepackage{babel}
  \usepackage[german=quotes]{csquotes}
  \usepackage{datetime2}
  \usepackage{ifthen}
  \usepackage{dingbat}
% ===========================================================

%% Fonts und Typographie
%% ===========================================================
  \usepackage{charter}
  \usepackage{nimbusmononarrow}
  \usepackage[babel=true,final,tracking=smallcaps]{microtype}
  \DisableLigatures{encoding = T1, family = tt* }
  \usepackage{ellipsis}
  \usepackage{listings}
  
% ===========================================================

% Farben
% ===========================================================
  \usepackage[usenames,x11names,table]{xcolor}
  \definecolor{fbblau}{HTML}{3078AB}
  \definecolor{blackberry}{rgb}{0.53, 0.0, 0.25}
% ===========================================================

% Mathe-Pakete und -Einstellungen
% ===========================================================
  \usepackage{mathtools}
  \usepackage{amssymb}
  \usepackage[bigdelims]{newtxmath}
  \allowdisplaybreaks
  \usepackage{bm}
  \usepackage{wasysym}
% ===========================================================

% TikZ
% ===========================================================
  \usepackage{tikz}
  \usepackage{tikz-cd}
  \usepackage{pgf}
  \usetikzlibrary{matrix}
  \usetikzlibrary{positioning}
  \usetikzlibrary{arrows}
  \usetikzlibrary{shapes}
  \tikzset{>=Latex}
  \usepackage[mode=image]{standalone}
% ===========================================================

% Seitenlayout, Kopf-/Fußzeile
% ===========================================================
  \usepackage{scrlayer-scrpage}
  \usepackage[top=3cm, bottom=2.5cm, left=2.5cm, right=2cm]{geometry}  
  \pagestyle{scrheadings}
  \clearscrheadfoot 
  \setheadsepline{0.4pt}
  \setfootsepline{0.4pt}
  \setkomafont{pagehead}{\normalfont\footnotesize}
  \setkomafont{pagefoot}{\normalfont\footnotesize}
  \lohead[]{Übung zur Vorlesung \textit{Embedded Systems Engineering} -- Wintersemester 2026}
  \rohead[]{%
    \ifthenelse{\boolean{loesungen}}{\textbf{Lösungen zu Blatt \blattnummer}}{Blatt \blattnummer}}
  %\rofoot[Version: \DTMnow]{%
  %  \ifthenelse{\boolean{loesungen}}{Version: \DTMnow}{}}
  %\lofoot[\jobname.tex]{%
  %  \ifthenelse{\boolean{loesungen}}{\jobname.tex}{}}
  \cfoot[\thepage]{\thepage}
  \raggedbottom
  \usepackage{setspace}          
  \linespread{1.2}
  \setlength{\parindent}{0pt}
  \setlength{\parskip}{0.5\baselineskip}
  \usepackage[all]{nowidow}
  \usepackage{lscape}
% ===========================================================

% Hyperref
% ===========================================================
  \usepackage[%
    hidelinks,
    pdfpagelabels,
    bookmarksopen=true,
    bookmarksnumbered=true,
    linkcolor=black,
    urlcolor=fbblau,
    plainpages=false,
    pagebackref,
    citecolor=black,
    hypertexnames=true,
    pdfborderstyle={/S/U},
    linkbordercolor=fbblau,
    colorlinks=false,
    backref=false]{hyperref}
  \hypersetup{final}
  \makeatletter
  \g@addto@macro{\UrlBreaks}{\UrlOrds}
  \makeatother
% ===========================================================

% Listen und Tabellen
% ===========================================================
  \usepackage{multicol}
  \usepackage{multirow}
  \usepackage[shortlabels]{enumitem}
  \setlist[enumerate]{font=\bfseries,itemsep=0.25\baselineskip}
  \setlist[itemize]{label=$\triangleright$,itemsep=-0.25\baselineskip}
  \usepackage{tabularx}
  \usepackage{array}
  \usepackage{listings}
  \lstset{breaklines=true,
    basicstyle=\ttfamily,
    frame=tblr,
    language=c,
    numbers=left}
  \usepackage{include/mips}

\definecolor{lightgreen}{rgb}{0,0.7,0}
\definecolor{gray}{rgb}{0.5,0.5,0.5}
\definecolor{mauve}{rgb}{0.58,0,0.82}

\lstdefinestyle{mips}{ %
  language=[mips]Assembler,       % the language of the code
  keywordstyle=\color{blue},          % keyword style
  keywordstyle=[2]\color{magenta},     % keyword style
  keywordstyle=[3]\color{red},          % keyword style
  commentstyle=\color{lightgreen},       % comment style
  stringstyle=\color{mauve},         % string literal style
  escapeinside={\%*}{*)},            % if you want to add a comment within your code
  morekeywords={*,...}               % if you want to add more keywords to the set
}

\newcommand{\mipsinline}[1]{\lstinline[style=mips,columns=fixed]{#1}}
\newcommand{\mipsregister}[1]{\mipsinline{\$#1}}
  \newcommand{\qed}{\hfill \ensuremath{\Box}}
  \newcommand{\cmd}[1]{\texttt{#1}}
% ===========================================================

% Sonstiges
% ===========================================================
  \usepackage[textsize=small,color=Red1!80!OrangeRed1!80]{todonotes}
  \usepackage{lipsum}
% ===========================================================

% Shortcuts and Math Commands
% ===========================================================
  \newcommand{\BB}{\mathbb{B}}      % Bool
  \newcommand{\CC}{\mathbb{C}}      % komplexe Zahlen
  \newcommand{\NN}{\mathbb{N}}      % nat. Zahlen
  \newcommand{\QQ}{\mathbb{Q}}      % rat. Zahlen
  \newcommand{\RR}{\mathbb{R}}      % reelle Zahlen
  \newcommand{\ZZ}{\mathbb{Z}}      % ganze Zahlen
  \newcommand{\oh}{\mathcal{O}}      % Landau-O
  \newcommand{\ol}[1]{\overline{#1}}    % überstreichen
  \newcommand{\wt}[1]{\widetilde{#1}}    % überschlängeln
  \newcommand{\wh}[1]{\widehat{#1}}    % überdachen
  \newcommand{\entspricht}{\mathrel{\widehat{=}}}
  
  \DeclarePairedDelimiter{\absolut}{\lvert}{\rvert}    % Betrag
  \DeclarePairedDelimiter{\ceiling}{\lceil}{\rceil}    % aufrunden
  \DeclarePairedDelimiter{\Floor}{\lfloor}{\rfloor}    % aufrunden
  \DeclarePairedDelimiter{\Norm}{\lVert}{\rVert}      % Norm
  \DeclarePairedDelimiter{\sprod}{\langle}{\rangle}    % spitze Klammern
  \DeclarePairedDelimiter{\enbrace}{(}{)}          % runde Klammern
  \DeclarePairedDelimiter{\benbrace}{\lbrack}{\rbrack}  % eckige Klammern
  \DeclarePairedDelimiter{\penbrace}{\{}{\}}        % geschweifte Klammern
  \newcommand{\Underbrace}[2]{{\underbrace{#1}_{#2}}}   % bessere Unterklammerungen
  
  % Kurzschreibweisen für Faule und Code-Vervollständigung
  \newcommand{\abs}[1]{\absolut*{#1}}
  \newcommand{\ceil}[1]{\ceiling*{#1}}
  \newcommand{\flo}[1]{\Floor*{#1}}
  \newcommand{\no}[1]{\Norm*{#1}}
  \newcommand{\sk}[1]{\sprod*{#1}}
  \newcommand{\enb}[1]{\enbrace*{#1}}
  \newcommand{\penb}[1]{\penbrace*{#1}}
  \newcommand{\benb}[1]{\benbrace*{#1}}
  \newcommand{\stack}[2]{\makebox[1cm][c]{$\stackrel{#1}{#2}$}}

% Übungszettel-Krams
% ===========================================================
  \usepackage[tikz]{mdframed}
  \newmdenv[%
    backgroundcolor = gray!20,
    linewidth=0pt,
    skipabove = 0.5cm,
    skipbelow = 0cm
  ]{notes}
  
  % Lösungsausgabe
  \newboolean{praesenzlsgn}
  \setboolean{praesenzlsgn}{\printpraesenzlsg}
  \newboolean{loesungen}
  \setboolean{loesungen}{\printloesungen}
  % Bewertungsschema-Ausgabe
  \newboolean{bewertungen}
  \setboolean{bewertungen}{\printbewertungen}

  \newcommand{\iforiginal}[1]{\ifthenelse{\boolean{praesenzlsgn}}{}{#1}}
  \newcommand{\ifpraesenzlsg}[1]{\ifthenelse{\boolean{praesenzlsgn}}{#1}{}}
  \newcommand{\ifloesung}[1]{\ifthenelse{\boolean{loesungen}}{#1}{}}
  \newcommand{\ifbewertung}[1]{\ifthenelse{\boolean{bewertungen}}{#1}{}}

  \newcommand{\umbruchoriginal}{\iforiginal{\newpage}}        % Seitenumbruch nur im Original
  \newcommand{\umbruchpraesenzlsg}{\ifpraesenzlsg{\newpage}}  % Seitenumbruch nur in Lösung d. Präsenzaufgabe
  \newcommand{\umbruchlsg}{\ifloesung{\newpage}}              % Seitenumbruch nur in Lösung
  \newcommand{\umbruchbewertung}{\ifbewertung{\newpage}}      % Seitenumbruch nur in Lösung mit Bewertungsschema
  \newcommand{\umbruchnurpraesenz}{\ifthenelse{\boolean{loesungen}}{}{\umbruchpraesenzlsg}}
  \newcommand{\umbruchnurlsg}{\ifthenelse{\boolean{bewertungen}}{}{\umbruchlsg}}

  \usepackage{comment}
  \newenvironment{praesenzlsg}%
    {\ifthenelse{\boolean{praesenzlsgn}}{\textbf{--- Präsenzaufgabe Anfang ---} \mbox{} \newline}{}}%
	{\ifthenelse{\boolean{praesenzlsgn}}{\mbox{} \newline \textbf{--- Präsenzaufgabe Ende ---}}{}}
  \newenvironment{loesung}%
    {\ifthenelse{\boolean{loesungen}}{\textbf{--- Lösung Anfang ---} \mbox{} \newline}{}}%
    {\ifthenelse{\boolean{loesungen}}{\mbox{} \newline \textbf{--- Lösung Ende ---} \ifthenelse{\boolean{bewertungen}}{}{\umbruchlsg}}{}}
  \newenvironment{bewertung}%
    {\ifthenelse{\boolean{bewertungen}}{\textbf{--- Bewertungsschema Anfang ---} \mbox{} \newline}{}}%
    {\ifthenelse{\boolean{bewertungen}}{\mbox{} \newline \textbf{--- Bewertungsschema Ende ---} \umbruchlsg}{}}

  \ifthenelse{\boolean{praesenzlsgn}}{}{\excludecomment{praesenzlsg}}
  \ifthenelse{\boolean{loesungen}}{}{\excludecomment{loesung}}
  \ifthenelse{\boolean{bewertungen}}{}{\excludecomment{bewertung}}

  % Aufgabennummerierung
  \newcounter{aufgabennummer} 
  \setcounter{aufgabennummer}{1}
  \newcommand{\aufgabe}[1]{\textbf{Aufgabe \blattnummer.\theaufgabennummer} \stepcounter{aufgabennummer} \hfill (#1 Punkte)}
  %\newcommand{\aufgabetitel}[2]{\iforiginal{\vspace{0.5cm}} \textbf{Aufgabe \blattnummer.\theaufgabennummer \ (#2)} \stepcounter{aufgabennummer} \hfill (#1 Punkte)}
  \newcommand{\aufgabetitel}[2]{\iforiginal{\vspace{0.5cm}} \textbf{Aufgabe \theaufgabennummer \ (#2)} \stepcounter{aufgabennummer} \hfill (#1 Punkte)}
  \newcommand{\praesenztitel}[1]{\iforiginal{\vspace{0.5cm}} \textbf{Aufgabe \blattnummer.\theaufgabennummer \ (#1)} \stepcounter{aufgabennummer} \hfill \smallpencil}

\newcommand{\danger}{\raisebox{0.5mm}{\fontencoding{U}\fontfamily{futs}\selectfont\char 66\relax}}
\newcommand{\ddanger}[1]{{\danger} \textbf{#1} {\danger}}

% Blatt-Thema (optional)
\newcommand{\blatthema}{} % default: leer
\newcommand{\setblatthema}[1]{\renewcommand{\blatthema}{#1}}

% Rechnerstrukturen-Krams
% ===========================================================
  \newcommand{\Band}{\ensuremath{\mathbin{\texttt{AND}}}}
  \newcommand{\Bor}{\ensuremath{\mathbin{\texttt{OR}}}}
  \newcommand{\Bnot}{\ensuremath{\mathbin{\texttt{NOT}}}}
  \newcommand{\Bnand}{\ensuremath{\mathbin{\texttt{NAND}}}}
  \newcommand{\Bnor}{\ensuremath{\mathbin{\texttt{NOR}}}}
  \newcommand{\Bxor}{\ensuremath{\mathbin{\texttt{XOR}}}}
  \newcommand{\Bxnor}{\ensuremath{\mathbin{\texttt{XNOR}}}}
  \newcommand{\Bone}{\ensuremath{\mathbin{\texttt{ONE}}}}
  \newcommand{\Bzero}{\ensuremath{\mathbin{\texttt{ZERO}}}}
  \newcommand{\Binary}[1]{\ensuremath{\texttt{#1}}\xspace}
  \newcommand{\ieee}[3]{\fbox{\ensuremath{\texttt{#1}}}\fbox{\ensuremath{\texttt{#2}}}\fbox{\ensuremath{\texttt{#3}}}}
  \DeclareMathOperator{\DNF}{DNF}
  \DeclareMathOperator{\KNF}{KNF}
  \DeclareMathOperator{\xor}{\oplus}
  \newcommand{\dc}{\star} %don't-care-Wert


% Create Code Listings
% ===========================================================
  \lstset{
  basicstyle=\ttfamily,
  numbers=left,
  numberstyle=\small,
  stepnumber=1,
  numbersep=10pt,
  frame=single,
  columns=fullflexible,
  keepspaces=true,
  showstringspaces=false,
  xleftmargin=1.5em
}
 

% Änderungen 2020: Hinweise auf Digitallehre angepasst; Rechnergeschichte entfernt; Blatt 1 und 2 zusammengefasst.
\begin{document}
\setblatthema{Embedded Systems Grundlagen}
\iforiginal{\thispagestyle{scrplain}
\vspace*{-3cm}
\begin{minipage}[t][1.1cm][c]{4.5cm}
  \includegraphics[width=4cm]{include/hwr-logo.png}
\end{minipage}
\hfill
\begin{minipage}[t][1.5cm][c]{8cm}
  \begin{center}
  \begin{footnotesize}
    \textsf{Prof. Dr.-Ing. Sebastian Schlesinger} \\[-0.1cm]
    \textsf{Fachbereich 2 - Duales Studium Wirtschaft \& Technik}
    %\url{https://www.hwr-berlin.de}
  \end{footnotesize}
  \end{center}
\end{minipage}
\hfill
%\begin{minipage}[t][1.6cm][c]{3.2cm}
%  \includegraphics[width=2.8cm]{include/ESlogo2}
%\end{minipage}

\vspace*{-0.3cm}
\begin{center} 
  \hrulefill \\[0.1cm]
  {\large Übung} \\[0.15cm]
    {\huge \bfseries Embedded Systems Engineering} \\[0.10cm]
  \ifthenelse{\equal{\blatthema}{}}{}{%
    {\Large \bfseries \blatthema} \\[0.10cm]
  }%
  {HWR Berlin, Wintersemester 2026} \\[-0.4cm]
  \begin{tabular}{lcr}
    \hspace{0.3\textwidth}   & \hspace{0.3\textwidth} & \hspace{0.3\textwidth} \\
    Prof. Dr.-Ing. Sebastian Schlesinger   &                        & "Ubung    \\ 
    Blatt \blattnummer     % \abgabetermin          \\ 
  \end{tabular} \\[0.1cm]
  \hrulefill
\end{center}}

% \begin{notes} \small
% 	\textbf{Abgabetermine für Blatt 1:}
	
% 	Aufgaben 1.2/1.3: Montag, 11. April, 09:00 Uhr \\
% 	Aufgabe 1.4: Mittwoch, 20. April, 18:00 Uhr
% \end{notes}
\aufgabetitel{$3$}{Embedded Systems - Grundlagen}\\
\begin{enumerate}
    \item Was sind ’eingebettete’ Systeme? Überlegen Sie sich Beispiele und klassifizieren Sie diese nach den in der Vorlesung vorgestellten Kriterien zur Klassifikation eingebetteter Software.
    \item Wie lässt sich der Aufbau eingebetteter Systeme beschreiben?
    \item Was bezeichnet der Begriff Regelstrecke?
\end{enumerate}

\begin{loesung}
\textbf{Zu 1: Begriff und Klassifikation}

\emph{Definition.} Ein eingebettetes System ist ein informationsverarbeitendes System aus Hard- und
Software, das in ein umgebendes technisches System integriert ist und dieses \emph{steuert bzw.\ regelt}.
Charakteristisch sind:
\begin{itemize}
    \item Es steuert eine (meist analoge) technische Umgebung mit digitalen Komponenten.
    \item Die Kopplung an die Umgebung erfolgt über \textbf{Sensoren} (Messgrößen) und \textbf{Aktoren} (Stellgrößen).
    \item Es ist für die Nutzerin in der Regel nicht als ,,Computer`` sichtbar und nicht frei programmierbar.
    \item Sein Zweck ist nicht die Berechnung einer mathematischen Funktion, sondern die \emph{dauerhafte
          Interaktion} mit der Umgebung; es terminiert im Normalbetrieb nicht.
    \item Es unterliegt typischerweise Ressourcen-, Energie-, Bauraum- und Stückkostenbeschränkungen.
\end{itemize}
Zur Einordnung: Über 98\,\% aller weltweit gefertigten Prozessoren stecken in eingebetteten Systemen.
Werden eingebettete Systeme vernetzt und mit der Cyber-Welt rückgekoppelt, spricht man von
\emph{Cyber-Physical Systems} bzw.\ \emph{Internet of Things}.

\emph{Klassifikationskriterien aus der Vorlesung:}
\begin{enumerate}[label=(\alph*)]
    \item transformatorisch vs.\ reaktiv
    \item diskretes vs.\ kontinuierliches Zeitverhalten
    \item verteilte vs.\ monolithische Systeme
    \item sicherheitskritisch vs.\ nicht sicherheitskritisch (z.\,B.\ Infotainment)
    \item Massenmarkt-Gesetzen unterliegend (Stückkosten-/Time-to-Market-Druck) vs.\ Kleinserie
\end{enumerate}
Ergänzend die in der Vorlesung gezeigte Schachtelung: \emph{eingebettete Systeme} $\supseteq$
\emph{reaktive Systeme} $\supseteq$ \emph{Echtzeitsysteme}, sowie die Unterscheidung \emph{harte} vs.\
\emph{weiche} Echtzeit.

\begin{table}[h]
\centering
\footnotesize
\renewcommand{\arraystretch}{1.25}
\resizebox{\textwidth}{!}{%
\begin{tabular}{|p{2.9cm}|p{2.2cm}|p{3.1cm}|p{2.8cm}|p{3.6cm}|p{2.6cm}|}
\hline
\textbf{Beispiel} & \textbf{transf./ reaktiv} & \textbf{Zeitverhalten} &
\textbf{Architektur} & \textbf{Kritikalität / Echtzeit} & \textbf{Markt} \\ \hline

ABS-/ESP-Steuer\-gerät &
reaktiv &
hybrid: kontinuierliche Radschlupf-Regelung, diskreter Regelzyklus &
verteilt (mehrere Steuergeräte, CAN/FlexRay) &
sicherheitskritisch, harte Echtzeit (ISO 26262) &
Massenmarkt \\ \hline

Herzschrittmacher &
reaktiv &
hybrid: kontinuierliche Messung, diskrete Stimulation &
monolithisch &
höchst sicherheitskritisch, harte Echtzeit &
Kleinserie (hoher Zertifizierungsaufwand) \\ \hline

Waschmaschinen\-steuerung &
überwiegend reaktiv; Programmablauf transformatorisch geprägt &
diskrete Ablaufsteuerung, kontinuierlich bei Temperatur/Drehzahl &
monolithisch &
kaum kritisch (Überlauf-/Verbrühschutz), weiche Echtzeit &
Massenmarkt \\ \hline

Infotainment-/ Navigationssystem &
reaktiv (Benutzerinteraktion) &
diskret &
verteilt (Fahrzeugbus, Mobilfunk/Cloud) &
nicht sicherheitskritisch, weiche Echtzeit; Rückwirkungsfreiheit nötig &
Massenmarkt \\ \hline

Industrieroboter &
reaktiv &
hybrid: kontinuierliche Achsregelung, diskrete Ablauflogik &
verteilt (SPS, Feldbus, Antriebe) &
sicherheitskritisch (Personenschutz), harte Echtzeit &
Kleinserie \\ \hline

Flugregler einer Drohne &
reaktiv &
kontinuierlich (Lageregelung im kHz-Bereich) &
monolithisch, bei Redundanz verteilt &
sicherheitskritisch, harte Echtzeit &
Massenmarkt (Consumer) \\ \hline

Smart Meter &
eher transformatorisch (Erfassen, Aggregieren, Übertragen) &
diskret &
verteilt &
nicht safety-, aber \emph{security}-kritisch; weiche Echtzeit &
Massenmarkt \\ \hline
\end{tabular}}
\end{table}

\emph{Grenzfälle (Diskussion aus der Vorlesung).} Smartphone, Smart-TV und Cloud-Dienst sind keine
eingebetteten Systeme im engeren Sinn: Sie sind frei programmierbare Allzweckrechner und steuern
keine technische Umgebung. Sie \emph{enthalten} allerdings eingebettete Teilsysteme (Baseband-Prozessor,
Kamera- und Energiemanagement-Controller). Eine Smart-Home-Anwendung ist ein Grenzfall: Die
Feldgeräte (Thermostat, Rollladenaktor) sind eingebettet, die App-/Cloud-Schicht nicht.

\textbf{Zu 2: Aufbau eingebetteter Systeme}

Der Aufbau lässt sich auf drei Ebenen beschreiben:

\emph{(a) Funktionale Struktur / Wirkungskette.} Die Kontrolleinheit greift über Aktoren in die
Umgebung ein und beobachtet sie über Sensoren -- es entsteht ein geschlossener Wirkungskreis.
Zusätzlich gibt es eine Benutzerschnittstelle zum Menschen.

\begin{center}
\begin{tikzpicture}[>=Latex, font=\small,
    box/.style={draw, rounded corners, minimum width=3cm, minimum height=0.9cm, align=center}]
  \node[box] (ctrl) {Kontrolleinheit\\(Prozessor + Software)};
  \node[box, above=1.1cm of ctrl] (ui) {Benutzerschnittstelle};
  \node[above=0.75cm of ui] (user) {Benutzer};
  \node[box, below left=1.7cm and 0.9cm of ctrl] (akt) {Aktoren};
  \node[box, below right=1.7cm and 0.9cm of ctrl] (sens) {Sensoren};
  \node[box, below=3.5cm of ctrl, minimum width=7cm] (strecke)
       {Regelstrecke / technische Umgebung};
  \draw[<->] (user) -- (ui);
  \draw[<->] (ui) -- (ctrl);
  \draw[->] (ctrl) -- node[above left, align=center, font=\footnotesize] {Stellgrößen} (akt);
  \draw[->] (akt) -- node[below left, align=center, font=\footnotesize] {Eingriff} (strecke);
  \draw[->] (strecke) -- node[below right, align=center, font=\footnotesize] {Messung} (sens);
  \draw[->] (sens) -- node[above right, align=center, font=\footnotesize] {Messgrößen} (ctrl);
\end{tikzpicture}
\end{center}

\emph{(b) Hardware-Architektur.} Ein oder mehrere Prozessoren mit zugehöriger Software, Speicher,
spezielle Hardware (ASIC, FPGA, Signalprozessoren) und Standardkomponenten; die Kopplung zur
technischen Umgebung erfolgt über Sensorik und Aktorik, die Kopplung untereinander über Busse
und Netzwerke (CAN, FlexRay, Ethernet).

\emph{(c) Software-Architektur.} Typischerweise geschichtet: Hardware-Abstraktion und Treiber,
darüber ein (Echtzeit-)Betriebssystem oder ein reiner Bare-Metal-Scheduler, Kommunikationsstacks
sowie die eigentliche Steuerungs-, Regelungs- und Diagnoselogik, realisiert als Tasks und
Interrupt-Service-Routinen.

\textbf{Zu 3: Regelstrecke}

Der Begriff stammt aus der Regelungstechnik. Die \textbf{Regelstrecke} (engl.\ \emph{plant},
\emph{controlled system}) ist derjenige Teil der technischen bzw.\ physikalischen Umgebung, dessen
Verhalten geregelt werden soll -- also das \emph{zu beeinflussende System} zwischen Stellgröße
(Eingang) und Regelgröße (Ausgang). Sie gehört ausdrücklich \emph{nicht} zum Rechner, sondern zum
Prozess.

Im Regelkreis unterscheidet man:
\begin{itemize}
    \item \textbf{Führungsgröße} $w$ (Sollwert), \textbf{Regelgröße} $x$ (Istwert),
          \textbf{Regeldifferenz} $e = w - x$
    \item \textbf{Stellgröße} $u$: Ausgang des Reglers, wirkt über den Aktor auf die Strecke
    \item \textbf{Störgröße} $z$: von außen auf die Strecke wirkende Einflüsse
\end{itemize}
Abzugrenzen ist die \emph{Regelung} (closed loop, mit Rückführung der Regelgröße) von der
\emph{Steuerung} (open loop, ohne Rückführung).

\emph{Beispiel Tempomat:} Regelstrecke ist die Längsdynamik des Fahrzeugs (Motor, Antriebsstrang,
Fahrwiderstände); Stellgröße ist das angeforderte Antriebsmoment, Regelgröße die Geschwindigkeit,
Störgrößen sind Steigung und Wind.

\emph{Bedeutung für eingebettete Systeme:} Die Regelstrecke verhält sich in aller Regel
\emph{kontinuierlich} (beschrieben durch Differentialgleichungen), der Regler arbeitet
\emph{zeitdiskret} mit fester Abtastzeit -- zusammen ergibt sich ein \emph{hybrides} System.
Abtastrate, Totzeiten und Jitter gehen unmittelbar in Regelgüte und Stabilität ein. Genau deshalb
ist Zeit in eingebetteten Systemen eine Frage der Korrektheit und nicht der Performanz.
\end{loesung}

\aufgabetitel{$3$}{Embedded Intro - Teil 1}\\
In \cite{lee2002embedded} steht ein Überblicksartikel zum Thema Embedded Software. Dieser gibt einen sehr guten
Einblick in die speziellen Probleme, denen die Entwickler eingebetteter Software gegenüberstehen. Lesen Sie die Abschnitte 1 bis 3 und bereiten Sie sich auf folgende Fragen
vor:
\begin{enumerate}
    \item Inwiefern unterscheidet sich Embedded Software von ’klassischer’ Software? Welche unterschiedlichen Anforderungen gibt es? Warum ist das für uns als Software-Entwickler wichtig und wie können wir damit umgehen?
    \item Diskutieren Sie die Begriffe Timeliness, Concurrency, Liveness, Interfaces, Heterogeneity, Reactivity und deren Bedeutung für die Entwickler eingebetteter Systeme.
    \item Welche Schwächen (laut Autor) haben die vorherrschenden Ansötze der Softwaretechnik für Embedded Software (Objektorientierung, Hardware-Beschreibungssprachen und Echtzeitbetriebssysteme)? Teilen Sie die Meinung des Autors?
\end{enumerate}

\begin{loesung}
\textbf{Zu 1: Embedded Software vs.\ klassische Software}

Lees Kernthese lautet: \emph{,,Embedded software is not just software on small computers.``} Der
Unterschied ist kein gradueller (weniger Speicher, langsamerer Prozessor), sondern ein prinzipieller.

\emph{Klassisches Bild von Software:} Ein Programm implementiert eine berechenbare Funktion
Eingabe~$\mapsto$~Ausgabe. Terminierung ist erwünscht, die Umgebung wird wegabstrahiert, und Zeit ist
eine reine Performanzfrage -- schneller ist besser. Die gesamte Abstraktionshierarchie der Informatik
(Befehlssatz, Programmiersprache, Betriebssystem) ist genau darauf ausgelegt, physikalische
Eigenschaften \emph{auszublenden}.

\emph{Embedded Software} dagegen:
\begin{itemize}
    \item ist in einen physikalischen Prozess eingebettet; die Wirkung in der physikalischen Welt
          \emph{ist} der Zweck und kein Seiteneffekt;
    \item terminiert nicht -- Nichttermination ist Anforderung, im klassischen Berechenbarkeitsmodell
          dagegen ein Fehlerfall;
    \item ist \emph{zeitbehaftet}: Korrektheit hängt vom Zeitpunkt ab. Häufig gilt nicht ,,so schnell
          wie möglich``, sondern ,,zum richtigen Zeitpunkt`` -- zu früh kann ebenso falsch sein wie
          zu spät;
    \item ist \emph{inhärent nebenläufig}, weil die physikalische Welt nebenläufig ist und nicht
          sequenziell;
    \item unterliegt harten Ressourcen-, Energie- und Stückkostenschranken;
    \item ist heterogen (Hardware/Software, kontinuierlich/diskret, verschiedene Protokolle);
    \item muss hochzuverlässig sein: Feldupdates sind teuer oder unmöglich, Fehler haben
          physikalische Konsequenzen.
\end{itemize}
Verschärfend kommt hinzu, dass moderne Rechnerarchitekturen (Caches, Pipelines, spekulative
Ausführung, Sprungvorhersage, dynamische Taktung) den \emph{Durchschnitt} verbessern und die
\emph{Vorhersagbarkeit} verschlechtern. Der Entwicklerin wird damit die Kontrolle über das
Zeitverhalten genommen -- genau über die Eigenschaft, die hier korrektheitsrelevant ist.

\emph{Warum ist das für uns wichtig, und wie gehen wir damit um?}
\begin{itemize}
    \item Nicht-funktionale Anforderungen (Zeit, Nebenläufigkeit, Ressourcen, Sicherheit) gehören in
          die \emph{Spezifikation} und in das \emph{Modell} -- sie lassen sich nicht nachträglich
          ,,hinein-optimieren``.
    \item Modellbasierte Entwicklung mit Modellen, die eine \emph{formale Semantik} besitzen:
          endliche Automaten, Statecharts, Timed Automata, signalflussorientierte Modelle
          (Simulink). Darauf aufbauend Simulation, Model Checking und Codegenerierung.
    \item Analysierbare Architekturen: statische Schedules, WCET-Analyse, Verzicht auf schwer
          vorhersagbare Konstrukte (dynamische Speicherallokation, unbeschränkte Rekursion,
          Garbage Collection), Coding Standards wie MISRA-C, speichersichere Sprachen wie Rust.
    \item Trennung von Funktionalität und Sicherheitsüberwachung: ein eigener Safety-Monitor bzw.\
          Watchdog neben der eigentlichen Steuerungslogik.
\end{itemize}

\textbf{Zu 2: Die Begriffe im Einzelnen}

\begin{itemize}
    \item \textbf{Timeliness (Zeitbehaftung, Rechtzeitigkeit).} Die Korrektheit hängt vom Zeitpunkt
          der Reaktion ab. Klassische Programmiersprachen und ihre Semantiken kennen keinen
          Zeitbegriff; auch klassische Temporallogik beschreibt nur \emph{qualitative} Zeit
          (,,immer``, ,,irgendwann``), aber keine exakten Schranken. \emph{Für Entwickler:}
          Deadlines, Perioden und Jitter müssen Teil von Spezifikation \emph{und} Analyse werden
          -- über Timed Automata, zeitbehaftete Logiken (TCTL), WCET- und Schedulability-Analyse.
    \item \textbf{Concurrency (Nebenläufigkeit).} Mehrere unabhängige Kontrollflüsse: Das System
          reagiert auf eine Menge von Sensoren und steuert ,,gleichzeitig`` Aktoren an. Die
          etablierten Lösungen der Betriebssystemwelt (Scheduling, Semaphore, Monitore, Interrupts)
          funktionieren, sind aber schwer vorhersagbar und fehleranfällig -- Race Conditions,
          Deadlocks, Prioritätsinversion. \emph{Für Entwickler:} nebenläufige Modelle mit klarer
          Semantik verwenden, Zustandsteilung minimieren (Message Passing statt Shared Memory),
          kritische Abschnitte verifizieren statt testen.
    \item \textbf{Liveness (Lebendigkeit).} Das System darf weder terminieren noch blockieren; es
          muss ,,immer wieder etwas Gutes passieren``. Auszuschließen sind Deadlock (alle Prozesse
          warten auf ein Ereignis, das nur ein anderer Prozess der Menge auslösen kann), Livelock
          und Starvation. \emph{Für Entwickler:} Lebendigkeitseigenschaften explizit formulieren
          (z.\,B.\ in LTL als $\Box\Diamond\,\varphi$) und per Model Checking nachweisen.
    \item \textbf{Interfaces (Schnittstellen).} In klassischer Software beschreibt eine Schnittstelle
          nur Typen von Parametern und Rückgabewerten und impliziert eine Übergabe des Kontrollflusses.
          Für eingebettete Systeme müsste sie darüber hinaus die \emph{Dynamik} beschreiben:
          zulässige Aufrufreihenfolgen und Protokolle, Zeitverhalten, Nebenläufigkeit,
          Ressourcenverbrauch. \emph{Für Entwickler:} Lee fordert ,,system-level types`` bzw.\
          Interface-Theorien, mit denen sich die Verträglichkeit einer Komposition ähnlich statisch
          prüfen lässt wie heute Datentypen.
    \item \textbf{Heterogeneity (Heterogenität).} Gemischt werden Hardware und Software,
          kontinuierliche und diskrete Anteile, synchrone und asynchrone Abläufe, blockierende und
          nichtblockierende Kommunikation, kontroll- und datenflussorientierte Sprachen.
          \emph{Für Entwickler:} Ein einziger Formalismus reicht nicht; erforderlich sind
          hierarchisch heterogene Modelle, eine sauber definierte Semantik an den Modellgrenzen und
          Co-Simulation.
    \item \textbf{Reactivity (Reaktivität).} Das System steht in kontinuierlicher Interaktion mit
          der Umgebung und arbeitet in \emph{deren} Tempo, nicht im eigenen. Ein ,,verpasstes``
          Ereignis kann schwere Folgen haben. \emph{Für Entwickler:} ereignisorientierte Modelle
          (Automaten, Statecharts), reaktive Sprachen (Esterel, Lustre, Ada) oder Erweiterungen
          etablierter Sprachen (SystemC, POSIX für C, RTSJ für Java); auf Implementierungsebene
          Interrupt-Architektur, Pufferung und definiertes Überlastverhalten.
\end{itemize}

\textbf{Zu 3: Schwächen der vorherrschenden Ansätze}

\emph{Objektorientierung.}
\begin{itemize}
    \item Sie abstrahiert genau das weg, worauf es ankommt: Zeit, Nebenläufigkeit und die
          physikalische Realisierung.
    \item Schnittstellen sind Methodensignaturen. Sie sagen nichts über Zeitverhalten, Protokoll
          oder Nebenläufigkeit aus; die Übergabe des Kontrollflusses (\emph{transfer of control})
          ist die falsche Abstraktion für nebenläufig interagierende Komponenten.
    \item Objekte sind passiv -- sie müssen aufgerufen werden. Nebenläufigkeit wird nachträglich
          über Threads und Monitore angeflanscht, ohne Kompositionssemantik: Aus der Korrektheit
          der Teile folgt nicht die Korrektheit der Komposition.
    \item Dynamische Bindung, dynamische Allokation und Garbage Collection erschweren die Vorhersage
          von Zeit- und Speicherverhalten.
\end{itemize}

\emph{Hardware-Beschreibungssprachen (VHDL, Verilog).}
\begin{itemize}
    \item Stärke: Zeit und Nebenläufigkeit sind Teil der Semantik (Discrete-Event-Modell) -- genau
          das, was der Objektorientierung fehlt.
    \item Schwäche: Das Abstraktionsniveau ist für komplexe Softwarefunktionalität zu niedrig; es
          fehlen Strukturierungs- und Abstraktionsmittel, wie sie große Softwaresysteme benötigen.
    \item Die Semantik ist auf Simulation und Synthese von Hardware zugeschnitten und in Details
          (Delta-Zyklen, Ereignisse mit gleichem Zeitstempel) subtil und wenig intuitiv.
          SystemC ist der Versuch, dieses Modell in die Softwarewelt (C++) zu übertragen.
\end{itemize}

\emph{Echtzeitbetriebssysteme.}
\begin{itemize}
    \item Zeitanforderungen werden nur \emph{indirekt} über Prioritäten ausgedrückt. Priorität ist
          eine schwache, globale und nicht-kompositionale Codierung einer Deadline.
    \item Daraus folgen die klassischen Probleme: Prioritätsinversion (prominentes Beispiel: Mars
          Pathfinder), Deadlocks, Starvation.
    \item Es gibt keine Kompositionalität: Das Hinzufügen einer Task kann das Zeitverhalten
          bestehender Tasks zerstören. Analysen wie RMS oder EDF setzen starke Annahmen voraus
          (Periodizität, Unabhängigkeit, bekannte WCET), die in realen Systemen selten alle gelten.
    \item Threads und Locks sind für Entwickler notorisch schwer korrekt zu verwenden; der
          resultierende Nichtdeterminismus macht Fehler schlecht reproduzierbar und Tests wenig
          aussagekräftig.
    \item Grundsätzlich: Ein RTOS löst ein Ressourcen-Multiplexing-Problem, liefert aber kein
          Modell, in dem Zeit Bestandteil der \emph{Semantik} ist.
\end{itemize}

\emph{Eigene Stellungnahme (Diskussionsaufgabe, es gibt keine eindeutig richtige Antwort).} Die
Analyse ist mehr als zwanzig Jahre später weitgehend bestätigt: Modellbasierte Entwicklung mit
Simulink/Stateflow, SCADE und Statecharts ist in Automotive und Avionik Standard, Timed Automata
und Model Checking sind etablierte Werkzeuge, und mit Rust setzt sich eine Sprache durch, die
Data-Race-Freiheit zur Übersetzungszeit garantiert. Einschränkend ist festzuhalten: Objektorientierung
und Echtzeitbetriebssysteme sind in der Praxis außerordentlich erfolgreich, weil sie Wiederverwendung,
Werkzeugunterstützung und verfügbares Personal mitbringen. Lees Kritik trifft sie weniger als
\emph{Implementierungstechnik}, sondern als \emph{alleinige Entwurfsabstraktion}. Der pragmatische
Schluss lautet daher: auf Entwurfsebene mit Modellen arbeiten, die Zeit und Nebenläufigkeit
semantisch erfassen, und die klassischen Techniken darunter als Implementierungsschicht nutzen.
\end{loesung}

\aufgabetitel{$3$}{Embedded Intro - Teil 2}\\
Lesen Sie nun die verbleibenden Abschnitte des Papers \cite{lee2002embedded} und bereiten Sie sich auf folgende Fragen vor: 
\begin{enumerate}
    \item Welchen Ansatz schlägt der Autor vor? Inwiefern vermeidet er die beschriebenen Schwächen der anderen Ansätze? Was halten Sie von dem Ansatz?
    \item Was bedeutet der Begriff Model of Computation und welchem Zweck dient er?
    \item Im Artikel werden die Models of Computation der folgenden Konzepte diskutiert. Notieren Sie zu jedem der genannten Konzepte 1-2 Stichpunkte, die bezeichnend für das jeweilige Model of Computation sind.
    \begin{enumerate}
        \item Datenflussmodelle
        \item Time Triggered Architectures (TTA)
        \item Synchronous / Reactive (SR)
        \item Discrete Event models (DE)
        \item Process Networks
        \item Rendezvous
        \item Publish and Subscribe
        \item Continuous Time models (CT) 
        \item Finite State Machines
    \end{enumerate}
    \item Was sind Vor- und Nachteile der Verwendung allgemeiner oder spezialisierter Modelle zur Modellierung eingebetteter Systeme? Welche Vor- und Nachteile hat die Kombination mehrerer Modelle?
\end{enumerate}

\begin{loesung}
\textbf{Zu 1: Lees Ansatz}

Lee schlägt vor, von der \emph{objektorientierten} zur \emph{aktorenorientierten} Entwurfssicht
(\emph{actor-oriented design}) überzugehen und die Semantik der Komposition explizit zu machen:
\begin{itemize}
    \item \textbf{Aktoren statt Objekte.} Komponenten sind nebenläufig, besitzen \emph{Ports} und
          kommunizieren über Datenströme bzw.\ Signale. Die Schnittstelle ist nicht der
          Methodenaufruf mit Kontrollflussübergabe, sondern der Datenaustausch.
    \item \textbf{Explizites Model of Computation.} Wie Komponenten nebenläufig ausgeführt werden
          und wie sie kommunizieren, legt nicht die Komponente fest, sondern das umgebende
          Berechnungsmodell (in Ptolemy~II durch einen \emph{Director} repräsentiert). Die
          ,,Gesetze``, nach denen die Komposition arbeitet, sind damit Teil des Modells und nicht
          implizite Eigenschaft der Implementierung.
    \item \textbf{Hierarchische Heterogenität.} Verschiedene Berechnungsmodelle werden hierarchisch
          geschachtelt, sodass jedes Teilproblem im dafür passenden Modell beschrieben wird -- etwa
          kontinuierliche Zeit für die Regelstrecke, endliche Automaten für Betriebsmodi,
          Datenfluss für Signalverarbeitung, Discrete Event für die Buskommunikation.
    \item \textbf{Interface-Theorien / ,,system-level types``.} Schnittstellen beschreiben auch
          Dynamik, Zeit und Nebenläufigkeit, sodass sich unverträgliche Kompositionen statisch
          erkennen lassen -- analog zur Typprüfung, aber auf Systemebene.
\end{itemize}

\emph{Inwiefern vermeidet das die beschriebenen Schwächen?} Zeit und Nebenläufigkeit sind Bestandteil
der Semantik und nicht nachträglich angeflanscht (Gegenentwurf zur Objektorientierung). Das
Abstraktionsniveau bleibt softwaretauglich und skaliert auf komplexe Funktionalität (Gegenentwurf zu
den Hardware-Beschreibungssprachen). Das Zeitverhalten ergibt sich aus dem Modell statt aus
Prioritäten und Scheduler-Heuristiken, und durch bewusst \emph{eingeschränkte}, domänenspezifische
Modelle wird Analysierbarkeit überhaupt erst möglich (Gegenentwurf zum RTOS-Ansatz).

\emph{Einschätzung (Diskussion).} Für den Ansatz spricht, dass er sich in der Praxis durchgesetzt hat,
wenn auch nicht unter Lees Namen: Simulink/Stateflow, SCADE, Modelica und AUTOSAR sind im Kern
aktorenorientierte, modellbasierte Werkzeugketten mit Codegenerierung; ROS\,2 realisiert ein
Publish-Subscribe-Berechnungsmodell. Kritisch anzumerken sind die starke Werkzeugabhängigkeit und
das Risiko des Vendor-Lock-in, der Einarbeitungsaufwand, die Gefahr auseinanderlaufender Modelle und
Implementierungen sowie der schwierigste Punkt des Ansatzes selbst: eine saubere Semantik \emph{an
den Grenzen} zwischen unterschiedlichen Berechnungsmodellen. Ptolemy~II selbst ist ein
Forschungswerkzeug geblieben -- die Ideen sind es nicht.

\textbf{Zu 2: Model of Computation}

Ein \emph{Model of Computation} (MoC, Berechnungsmodell) ist die Menge der Regeln, die die
\emph{Interaktion von Komponenten} innerhalb eines Modells bestimmen. Lee beschreibt es bildhaft als
die ,,Gesetze der Physik``, denen die Komponenten einer Komposition unterliegen. Ein MoC legt fest:
\begin{itemize}
    \item Was ist eine Komponente, und wann wird sie aktiv (Aktivierungs-/Feuerregel, Scheduling)?
    \item Wie wird kommuniziert -- blockierend oder nichtblockierend, gepuffert oder synchron,
          gerichtet oder anonym?
    \item Welche Rolle spielt Zeit: gar keine, logische Ordnung, diskrete Ticks, Zeitstempel oder
          kontinuierliche Zeit? Was bedeutet ,,gleichzeitig``?
    \item Wann ist eine Komposition wohldefiniert, determiniert und ausführbar?
\end{itemize}
\emph{Zweck:}
\begin{itemize}
    \item \textbf{Formale Semantik} als Grundlage für Simulation, Verifikation und Codegenerierung --
          ohne sie ist ein Modell nur ein Bild.
    \item \textbf{Kompositionalität:} Das Verhalten der Zusammensetzung folgt aus den Regeln und muss
          nicht für jede Kombination neu ermittelt werden.
    \item \textbf{Analysierbarkeit durch gezielte Einschränkung} der Ausdrucksstärke -- etwa
          statisches Scheduling, beschränkte Puffergrößen oder entscheidbare Deadlock-Freiheit.
    \item \textbf{Passung zur Domäne} und damit eine gemeinsame Sprache zwischen Domänenexpertin und
          Softwareentwickler.
\end{itemize}

\textbf{Zu 3: Charakteristika der einzelnen Berechnungsmodelle}

\begin{enumerate}[label=(\alph*)]
    \item \textbf{Datenflussmodelle.} Aktoren feuern datengetrieben, sobald genügend Token an ihren
          Eingängen anliegen; ein globaler Zeitbegriff fehlt. Im Spezialfall \emph{Synchronous
          Dataflow} (SDF) sind die Token-Raten statisch bekannt, wodurch sich Schedule, Puffergrößen
          und Deadlock-Freiheit bereits zur Übersetzungszeit bestimmen lassen. Natürlicher
          Anwendungsbereich: Signalverarbeitung und Streaming.
    \item \textbf{Time Triggered Architectures (TTA).} Rechnen und Kommunizieren folgen einem
          \emph{statischen, global zeitgesteuerten Fahrplan} (TDMA) auf Basis einer synchronisierten
          globalen Zeit; nichts geschieht ereignisgetrieben. Das ergibt maximale Vorhersagbarkeit,
          Testbarkeit und Fehlertoleranz, erkauft mit Unflexibilität gegenüber sporadischen
          Ereignissen und einer Auslegung sämtlicher Ressourcen auf den Worst Case.
    \item \textbf{Synchronous / Reactive (SR).} Es gilt die \emph{Synchronie-Hypothese}: Reaktionen
          benötigen keine Zeit, das System durchläuft diskrete Ticks, in denen alle Komponenten
          gleichzeitig auf die anliegenden Eingaben reagieren. Die Semantik ist ein Fixpunkt, weshalb
          eine Kausalitätsanalyse nötig ist (Ausschluss von Kausalitätsschleifen). Determiniert und
          sehr gut analysierbar (Esterel, Lustre, SCADE), aber global eng gekoppelt und damit auf
          verteilten Plattformen teuer.
    \item \textbf{Discrete Event models (DE).} Ereignisse tragen \emph{Zeitstempel}; die Ausführung
          verarbeitet sie in chronologischer Reihenfolge aus einer globalen Ereigniswarteschlange.
          Zeit ist explizit Teil der Semantik. Grundlage von Hardware- und Netzwerksimulatoren
          (VHDL, Verilog, SystemC). Kritische Punkte: gleichzeitige Ereignisse und
          Null-Verzögerungs-Rückkopplungen (Delta-Zyklen, ,,superdense time``) sowie die teure
          globale Zeitordnung in verteilten Systemen.
    \item \textbf{Process Networks (Kahn Process Networks).} Nebenläufige Prozesse kommunizieren über
          konzeptuell unbeschränkte FIFO-Kanäle; Lesen blockiert, Schreiben blockiert nie. Das
          Ergebnis ist \emph{determiniert und unabhängig von der Ausführungsreihenfolge} -- eine sehr
          starke Eigenschaft. Ein Zeitbegriff fehlt; Beschränktheit der Puffer und Deadlock-Freiheit
          sind im Allgemeinen unentscheidbar.
    \item \textbf{Rendezvous.} Synchrone, ungepufferte Punkt-zu-Punkt-Kommunikation: Sender und
          Empfänger warten aufeinander, der Datenaustausch findet atomar im ,,Treffpunkt`` statt
          (CSP, CCS, Occam, Ada). Gut geeignet zur Modellierung geteilter Ressourcen und formal
          hervorragend analysierbar (Prozessalgebren, Bisimulation); dafür starke Kopplung und
          Deadlock-Gefahr, Nichtdeterminismus über bewachte Auswahl.
    \item \textbf{Publish and Subscribe.} Anonyme, lose Kopplung über benannte Ereignisströme
          (Topics); Produzenten kennen ihre Konsumenten nicht. Sehr gut für dynamische, verteilte und
          skalierbare Systeme (heute DDS/ROS\,2, MQTT). Der Preis ist eine schwach definierte
          Semantik bezüglich Reihenfolge, Zeit und Zustellgarantien -- Determinismus und harte
          Echtzeit sind entsprechend schwer zuzusichern.
    \item \textbf{Continuous Time models (CT).} Das Verhalten wird durch Differentialgleichungen über
          einem kontinuierlichen Zeitbereich beschrieben; ausgeführt wird es näherungsweise durch
          einen numerischen Solver (Schrittweitensteuerung, Nullstellenerkennung). Es ist das
          passende Modell für die physikalische Regelstrecke; in Kombination mit diskreten Anteilen
          entstehen \emph{hybride} Systeme (Simulink, Modelica).
    \item \textbf{Finite State Machines.} Endliche Zustandsmenge mit Transitionen; das Verhalten ist
          vollständig aufzählbar und damit sehr gut verifizierbar (Erreichbarkeitsanalyse, Model
          Checking). Für sich genommen kennen FSM weder Nebenläufigkeit noch Zeit, und bei
          Komposition droht Zustandsexplosion -- daher die Erweiterungen: Statecharts (Hierarchie und
          Nebenläufigkeit), Timed Automata (Zeit) und hybride Automaten (FSM kombiniert mit CT).
\end{enumerate}

\textbf{Zu 4: Allgemeine vs.\ spezialisierte Modelle}

\emph{Allgemeine Modelle} (Turing-mächtige Sprachen, UML, sehr ausdrucksstarke Berechnungsmodelle):
\begin{itemize}
    \item \emph{Vorteile:} Alles ist modellierbar, ein Formalismus deckt das Gesamtsystem ab, keine
          Modellierungs-Verrenkungen nötig; breite Werkzeug- und Personalbasis.
    \item \emph{Nachteile:} Analyseaussagen werden schwach oder unentscheidbar (Halteproblem,
          Erreichbarkeit); statische Garantien zu Zeit und Speicher sind nicht zu bekommen; die
          Semantik ist häufig unterspezifiziert -- bei UML bedeutet dasselbe Diagramm in
          verschiedenen Werkzeugen Verschiedenes. Zudem sind solche Modelle selten domänennah.
\end{itemize}

\emph{Spezialisierte Modelle} (SDF, SR, Timed Automata, TTA, CT):
\begin{itemize}
    \item \emph{Vorteile:} Die eingeschränkte Ausdrucksstärke \emph{kauft} Analysierbarkeit --
          statisches Scheduling, beschränkter Speicherbedarf, Determinismus, Deadlock-Freiheit,
          Model Checking, WCET-Analyse und verlässliche Codegenerierung. Die Notation ist nah an der
          Domäne und damit für Domänenexperten lesbar.
    \item \emph{Nachteile:} Nicht alles lässt sich ausdrücken. Erzwungene Workarounds zerstören
          häufig genau die Analysevorteile, für die man das Modell gewählt hat. Hinzu kommen
          Lernaufwand, Werkzeugvielfalt und Lock-in-Risiken; die Grenzen des Modells müssen verstanden
          sein.
\end{itemize}

\emph{Kombination mehrerer Modelle (hierarchische Heterogenität):}
\begin{itemize}
    \item \emph{Vorteile:} Jedes Teilproblem wird im passenden Modell beschrieben und behält dort
          seine maximale Analysierbarkeit. Das entspricht der tatsächlich heterogenen Natur
          eingebetteter Systeme -- kontinuierliche Regelstrecke, diskrete Steuerung,
          ereignisbasierte Kommunikation. Lokale Verifikation und Wiederverwendung von Komponenten
          werden dadurch erst praktikabel.
    \item \emph{Nachteile:} Die Semantik \emph{an den Modellgrenzen} muss präzise definiert sein --
          etwa die Abtastung beim Übergang CT~$\to$~FSM oder die Zeitstempel-Interpretation beim
          Übergang DE~$\to$~SR. Genau dort entstehen erfahrungsgemäß die Fehler. Globale
          Eigenschaften wie Ende-zu-Ende-Deadlines lassen sich über Modellgrenzen hinweg deutlich
          schwerer beweisen als innerhalb eines Modells. Hinzu kommen konzeptionelle Komplexität,
          Integration der Werkzeugkette, Co-Simulation und Schulungsaufwand.
\end{itemize}

Lees Antwort lautet entsprechend: Die Kombination ist \emph{notwendig}, aber nur mit wohldefinierten
Schnittstellen zwischen den Berechnungsmodellen. In dieser Vorlesung sehen wir genau dieses Muster
am Robotikbeispiel: Statecharts für die Betriebsmodi, Timed Automata für Deadlines und Watchdogs,
Simulink für den Regelkreis und ROS\,2 als Publish-Subscribe-Schicht für die Kommunikation.
\end{loesung}

\begin{thebibliography}{9}

\bibitem{lee2002embedded}
Edward A. Lee.
\newblock Embedded software.
\newblock In: {\em Advances in Computers}, Vol.~56, pp.~55--95.
\newblock Elsevier, 2002.
\newblock DOI: \url{https://doi.org/10.1016/S0065-2458(02)80004-3}

\end{thebibliography}

\end{document}
